\chapter{Continuous records without repeated long reset gaps}
\label{ch:ci-dependence}

A long blind interval before every measurement makes conditional means nearly
stationary, but can waste most of the physical evolution time. We replace that
acquisition rule by equally spaced blocks and use a calibrated mixing envelope
to control their dependence. The improvement concerns declared evolution time;
it does not make readout energy or mixing calibration free. The concentration
ingredient is established Markov-chain theory, combined here with a whole-record
martingale bound and the continuous reversible comparison class.

\section{An apparatus contract for a dependent record}
Let $X_t$ be a finite stationary continuous-time Markov chain with law $\pi$ and
generator $Q$. Write $Q^\dagger$ for its adjoint in $L^2(\pi)$. Suppose an
\emph{externally justified} lower bound $\lambda>0$ on the spectral gap of
$-(Q+Q^\dagger)/2$ is available. No complete generator need be supplied.
The bound is required for every model to which the confidence claim applies;
fitting it to the same coarse record does not establish this contract.

Divide acquisition into deterministic blocks of duration $H$. The hidden
block starts $X_i=X_{(i-1)H}$ have transition $P=e^{HQ}$. In block $i$, a fixed
setting produces a statistic $Z_i\in[0,1]$ from a noninvasive record contained
in that block. Unobserved padding is physical evolution and counts in $H$.
The record kernel depends on the current hidden start and its predetermined
setting, with fresh conditional detector noise. A detector with memory must
instead be included in the state for which the mixing guarantee holds.

For an agreement pair at lag $t_i\leq H$, two fresh detector readings may occur
at a shared boundary between blocks. Reusing the same noisy detector bit as
both the previous terminal and next initial reading generally violates the
contract. A deterministic state readout has no such extra random memory.
Settings, quotas and contrast coefficients are fixed before the inference
data, or using an independent pilot. The result below does not cover adaptive
routing of settings based on earlier outcomes.

\begin{theorem}[Whole-block concentration under calibrated contraction]
Define
\begin{equation}
 \rho=e^{-\lambda H},\qquad
 A_H=\frac{1+\rho}{1-\rho},\qquad
 K_H=(\sqrt{A_H}+1)^2.
 \label{eq:ci-dependent-proxy}
\end{equation}
For deterministic real weights $c_i$, let
$\mu_i=\mathbb E_\pi Z_i$. Then
\begin{equation}
 \mathbb E_\pi\exp\!\left(\theta\sum_i c_i(Z_i-\mu_i)\right)
 \leq\exp\!\left(\frac{\theta^2K_H}{8}\sum_i c_i^2\right).
 \label{eq:ci-dependent-mgf}
\end{equation}
Consequently, for $\alpha\in(0,1)$,
\begin{equation}
 \Pr_\pi\!\left\{\sum_i c_i(Z_i-\mu_i)>
 \sqrt{\frac{K_H\log(1/\alpha)}2\sum_i c_i^2}\right\}
 \leq\alpha.
 \label{eq:ci-dependent-tail}
\end{equation}
If the initial hidden law has total variation distance at most $\beta$ from
$\pi$, the same tail event has probability at most $\alpha+\beta$.
\end{theorem}
\begin{proof}
For centered $f$, dissipativity gives
\[
 \frac{d}{dt}\|e^{tQ}f\|_\pi^2
 =2\langle e^{tQ}f,\tfrac12(Q+Q^\dagger)e^{tQ}f\rangle_\pi
 \leq-2\lambda\|e^{tQ}f\|_\pi^2.
\]
Thus $\|P-\Pi\|_{L^2(\pi)}\leq\rho$, also for nonreversible $Q$.
Set $g_i(x)=\mathbb E[Z_i\mid X_i=x]$ and
$D_i=Z_i-g_i(X_i)$. Relative to the filtration containing the past blocks and
the current hidden start, $D_i$ is conditionally centered with range length
at most one. Conditional Hoeffding therefore bounds the moment generating
function of $\sum c_iD_i$ by $\exp(\theta^2\sum c_i^2/8)$.

The stationary absolute-contraction result of Fan, Jiang and Sun,
Theorem~1~\cite{ciFan2021}, applies to the deterministic, possibly
time-dependent state functions $c_i g_i$. It bounds the moment generating
function of $\sum c_i(g_i(X_i)-\pi g_i)$ by
$\exp(\theta^2 A_H\sum c_i^2/8)$. The two sums need not be independent.
Apply H\"older with conjugate exponents
$p=1+A_H^{-1/2}$ and $q=1+A_H^{1/2}$. The resulting coefficient is
$pA_H+q=(\sqrt{A_H}+1)^2$, proving the moment bound. Chernoff optimization
gives the one-sided tail.

Finally, the entire future observed path is a Markov kernel of the initial
state. Total variation contraction transfers initial distance $\beta$ to
the full path law. It is paid once, not once per record.
\end{proof}

This is a conservative bound for a statistic of an entire block. In particular,
$K_H\to4$ as $\rho\to0$; it does not recover the sharp independent-record
constant. Shared hidden endpoints can create dependence even with fresh
detector noise. Exact independent preparations retain their separate,
stronger concentration guarantee.

\section{From counts to physical inference}
For $m$ fixed settings with quotas $n_k$, probabilities $p_k$ and coefficients
$a_k$, apply the theorem with $c_i=a_k/n_k$ in setting $k$. If
$\widehat p_k=S_k/n_k$, a valid stationary contrast lower bound is
\begin{equation}
 L_a=\sum_k a_k\widehat p_k-
 \sqrt{\frac{K_H\log(1/\alpha)}2\sum_k\frac{a_k^2}{n_k}}
 -\sum_k|a_k|\delta_k,
 \label{eq:ci-dependent-contrast}
\end{equation}
where $\delta_k$ are separately valid whole-record calibration radii.
Subtracting the continuous reversible support $h_{\mathcal M}(a)$ gives
$g=[L_a-h_{\mathcal M}(a)]_+$. Under the stationary ordinary-time-reversal
and physical entropy assumptions of the preceding construction, this yields
\begin{equation}
 \sigma\geq\frac{8g^2}{(\sum_k|a_k|\sqrt{t_k})^2}.
 \label{eq:ci-dependent-entropy}
\end{equation}
Here $t_k$ is the observed path duration; the sampling cost uses the complete
block duration $H$, including padding. The mixing assumption licenses a
confidence calculation. It does not itself establish the thermodynamic
interpretation or local detailed balance.

Alternatively, a simultaneous probability box uses radii
$\sqrt{K_H\log(2m/\alpha)/(2n_k)}+\delta_k$. A score or contrast may be
selected after seeing that jointly valid box. The smaller fixed-contrast
radius in~\eqref{eq:ci-dependent-contrast} does not permit such reuse.
The code evaluates exponentials, logarithms and square roots with outward
rational bounds inherited from the certified arithmetic layer.

A single initial blind interval $G$ suffices if a stationary mass lower bound
$\pi_{\min}\geq p_*>0$ is available. The contraction estimate implies
\begin{equation}
 \sup_x\|\delta_xe^{GQ}-\pi\|_{\mathrm{TV}}
 \leq\min\left\{1,\frac12\sqrt{\frac{1-p_*}{p_*}}e^{-\lambda G}\right\}.
 \label{eq:ci-one-burn}
\end{equation}
This is an initial-law distance, not a record-by-record conditional bias and
not an observed event that a reset succeeded. A separate, externally valid
calibration failure allowance $\gamma$ increases total failure to
$\alpha+\beta+\gamma$. A postselected claim that calibration looked adequate
cannot replace this requirement.

\section{A frozen comparison at equal evolution time}
The initial evaluation fixes $\lambda=1/500$, $p_*=1/8$, lags $1,2$,
$a=(-1,1)$, and $H=2$. New acquisition pays $G=4000$ once, with
$\alpha=0.049$ and $\beta=0.001$. Its total failure allowance is $0.05$.
The old conditional-mean method pays a blind gap before every pair. To avoid
a weak baseline, its gap is optimized before any data over
$\{0,50,\ldots,10000\}$, minimizing its certified contrast halfwidth at the
same total evolution time. Both methods count padding and every blind gap.

The valid equilibrium model has four states $(s,b)$. The sector flips at
$10^{-3}$; the observed bit flips at rates $0.02$ and $2$ in the two sectors.
Its uniform stationary law is reversible and its additive gap is $0.002$.
The driven alternative is an eight-state ring with clockwise rate $4$,
counterclockwise rate $0.1$, and a detector separating the two half-rings.
Its true agreement contrast is $0.1469539766$. Both start at a fixed hidden
state before acquisition. A third equilibrium model has a sector rate
$10^{-5}$, violating the advertised gap by a factor of $100$; it starts
stationary to isolate the incorrect mixing calibration.

\begin{center}
\begin{tabular}{r|rr|rr|rr}
Evolution & \multicolumn{2}{c|}{Halfwidth} &
\multicolumn{2}{c|}{Endpoint readouts} &
\multicolumn{2}{c}{Driven exclusions /1000}\\
budget & New & Old & New & Old & New & Old\\\hline
$200000$ & $.18327$ & $.29032$ & $196000$ & $176$ & $0$ & $84$\\
$400000$ & $.12894$ & $.21176$ & $396000$ & $332$ & $1000$ & $212$\\
$800000$ & $.09094$ & $.15422$ & $796000$ & $592$ & $1000$ & $465$
\end{tabular}
\end{center}

The new method gains substantial power at the two larger budgets but does
not dominate every operating point: at the smallest budget its certified
threshold exceeds the true contrast and all runs fail to certify. On the
valid equilibrium model, new false exclusions are $0/1000$ at each budget;
the old counts are $1,0,1$. Treating exactly the same new dependent counts
as independent gives $271,217,131$ false exclusions. Correct stationary
marginals alone therefore do not justify the independent analysis.

The wrong-envelope control produces $173,256,258$ false exclusions under the
new calculation, and $97,131,132$ under the old one. Neither method repairs
an invalid external mixing guarantee. These controls are outside the theorem,
not failures of its stated guarantee.

Reported Monte Carlo intervals are approximate 95\% Wilson intervals,
separate from the rigorous statistical confidence guarantee. For example,
$1000/1000$ corresponds to $[0.99617,1]$ and $0/1000$ to $[0,0.00383]$.
The old driven powers at the larger budgets are $0.212$ and $0.465$, with
intervals $[0.18779,0.23841]$ and $[0.43428,0.49599]$.
These simulations are not physical experimental validation.

\section{Trading readouts against waiting}
The dense-block improvement requires many more readouts. A second evaluation,
frozen separately after the first and using independent seeds, varies $H$
while leaving the observed lags and confidence contract unchanged. For total
evolution budget $T$, equal quotas are
\[
 n=\left\lfloor\frac{T-G}{2H}\right\rfloor,
 \qquad r_H=\sqrt{K_H\log(1/\alpha)/n}.
\]
Increasing $H$ reduces the retained data and weakens this bound, but can leave
enough margin for a useful test. It also reduces the number of detector
operations. All settings in the frozen menu are retained, including
underpowered choices; no score is selected using its own inference data.

\begin{center}
\begin{tabular}{r|r|r|r|r}
Evolution budget & $H$ & Endpoint readouts & Halfwidth & Driven exclusions /1000\\\hline
$400000$ & $10$ & $79200$ & $.13576$ & $983$\\
$400000$ & $20$ & $39600$ & $.14088$ & $816$\\
$400000$ & $50$ & $15840$ & $.15107$ & $339$\\
$800000$ & $20$ & $79600$ & $.09937$ & $1000$\\
$800000$ & $50$ & $31840$ & $.10655$ & $1000$\\
$800000$ & $100$ & $15920$ & $.11472$ & $994$\\
$800000$ & $200$ & $7960$ & $.12656$ & $892$\\
$800000$ & $500$ & $3184$ & $.15210$ & $394$
\end{tabular}
\end{center}

At budget $400000$, $H=10$ retains $98.3\%$ observed power with five times
fewer endpoint readouts than the original $H=2$ acquisition. At budget
$800000$, $H=100$ retains $99.4\%$ with fifty times fewer readouts. Their
approximate Wilson intervals are $[0.97294,0.98936]$ and
$[0.98697,0.99725]$. Every valid equilibrium refinement row has zero false
exclusions in $1000$ trials. The wrong-mixing control still fails; for example,
$H=100$ at budget $800000$ produces $198/1000$ false exclusions. Sparse
acquisition reduces detector operations without eliminating the external
mixing assumption. Nor does it guarantee high power: the last row retains
only $39.4\%$ observed power.

For apparatus design, suppose a constant waiting power $P_{\rm wait}$ and
energy per endpoint readout $E_{\rm read}$ have been measured. Relative to
one lag-unit $t_0$, define
$\zeta=E_{\rm read}/(P_{\rm wait}t_0)$. In units of $P_{\rm wait}t_0$,
the accounting model is
\begin{equation}
 C=T+\zeta R,\qquad
 n=\left\lfloor\frac{C-G}{2(H+2\zeta)}\right\rfloor,
 \qquad R=4n.
 \label{eq:ci-readout-cost}
\end{equation}
The code minimizes the outward-certified radius over the predeclared finite
grid $H=2,3,\ldots,1000$. Without integer quotas the objective is proportional
to $(H+2\zeta)K_H$. This supplies a reproducible tradeoff once costs are
calibrated, not a globally optimal protocol or a measured energy advantage.
Detector duration, preparation overhead and external calibration must also
be included when they are material. The current simulation records those
quantities as unmeasured rather than zero.

For example, at normalized budget $800000$, the finite-grid choices for
$\zeta=0,1,10,100$ are respectively $H=2,16,74,307$, with certified
halfwidths $0.09094,0.10401,0.12489,0.17558$. The $\zeta=10$ plan uses
$16936$ endpoint readouts and evolution time $630632$, consuming normalized
cost $799992$. These are prospective design calculations, not additional
held-out power estimates at those settings.

\section{Why some mixing cost cannot be removed}
A simple experiment-independent obstruction complements the usable bound.
Consider two fully observed stationary two-state chains. Both have rate
$0\to1$ equal to $\kappa$; their rates $1\to0$ are $3\kappa$ and
$\kappa/3$. Their stationary probabilities of state $1$ are $1/4$ and $3/4$.
Even a detector retaining every transition for duration $T$ sees the identical
constant-zero path under both models with common mass
$\tfrac14 e^{-\kappa T}$, and the identical constant-one path with common
mass $\tfrac14 e^{-3\kappa T}$. Hence equal-prior classification error obeys
\begin{equation}
 \varepsilon^*(T)\geq\frac18
 \left(e^{-\kappa T}+e^{-3\kappa T}\right).
 \label{eq:ci-slow-mixing-lower}
\end{equation}
An estimator of stationary occupation with error strictly below $1/4$ under
both hypotheses would also classify them by thresholding at $1/2$. Thus
uniformly small failure requires time of order $1/\kappa$ in this class, even
with full continuous observation. This overlap argument does not prove the
present bound optimal, but rules out a universal removal of slow-mixing cost
by simply sampling one trajectory more frequently.

The contribution of this chapter is the integrated acquisition and inference
contract: whole-record dependence, one initial calibration payment, exact
continuous-class support, explicit physical time, and a readout-cost design
frontier. Markov concentration, conditional Hoeffding and the overlap argument
are established tools rather than new physical laws.
